% SPDX-License-Identifier: Apache-2.0
\section{An Input Held at a Constant}
\label{sec:x-tie}

A unit of units joins each of its children's ports to one of its own
ports, or to a channel or a wire it makes in \code{run}. Sometimes an
input has nothing to be joined to: the design wants it at one value
for good, a step of three or an enable that is always on.
\code{tie(v)} is that value, written where the input is passed.

In the run, \code{tie(v)} is a wire set to \code{v} once and never
driven again, so the child reads \code{v} on every step. In the
parent's netlist it is a wire of the parent's own, named for the
child and numbered, driven by the constant and joined to the child's
input. The constant is evaluated when the parent is lowered as well
as when it runs, so it is a literal or a value made from literals, and
the two cannot disagree.

Before \code{tie}, a design that wanted an input held made a
\code{signal}, drove nothing into it, and passed its reading end.
That holds the type's default, which is a constant, but always zero,
and nothing in the source said it was meant.

A constant drives what it is joined to, so it may only be joined to an
input. Joined to an output it would be driven twice, and a channel is
three nets and not one, so the lowering refuses both, naming the wire,
the port and the child.

The accumulator below adds its input and a step to a total while it
is enabled; the parent holds the step at three and the enable on. The
build checks the parent's netlist against this run under nvc and
Verilator.

\begin{figure}[htbp]
\centering
\input{tie_timing.tex}
\caption{The accumulator with its step tied to three and its enable
on: the total grows by the input and three on every edge, and wraps.}
\label{fig:x-tie}
\end{figure}

\lstinputlisting[caption={\code{ex\_tie.rs}. Run with \code{bazel run //lib/examples:ex\_tie}.},label={lst:x-tie}]{lib/examples/ex_tie.rs}

\lstinputlisting[style=ascii,caption={What \code{ex\_tie} printed when the build ran it: the total and the output cycle by cycle, then the Verilog, the two constants the wires \code{acc\_tie0} and \code{acc\_tie1}.},label={lst:x-tie-out}]{out_tie.txt}
